\par\Needspace{14\baselineskip}
\SourceLeaf{139}
\begin{center}\textit{FIRST SECTION.}\\[4pt]
\textit{Relations between contiguous functions:}\end{center}
\SixthArticle{7}
We call a \textit{function} \textit{contiguous} to $F(\alpha,\beta,\gamma,x)$ if it arises from that function by increasing or decreasing the first, second or third element by unity, while the other three elements remain unchanged. Thus the primary function $F(\alpha,\beta,\gamma,x)$ gives six contiguous functions, and between any two of these and the primary function itself there is a very simple linear equation. We set out these fifteen equations here, omitting for brevity the fourth element, which is always understood to be $=x$, and denoting the primary function simply by $F$.
\[
\begin{array}{rcl}
{[1]}&0=&(\gamma-2\alpha-(\beta-\alpha)x)F+\alpha(1-x)F(\alpha+1,\beta,\gamma)\\
&&{}-(\gamma-\alpha)F(\alpha-1,\beta,\gamma)
\end{array}
\]
\[
\begin{array}{rcl}
{[2]}&0=&(\beta-\alpha)F+\alpha F(\alpha+1,\beta,\gamma)-\beta F(\alpha,\beta+1,\gamma)
\end{array}
\]
\[
\begin{array}{rcl}
{[3]}&0=&(\gamma-\alpha-\beta)F+\alpha(1-x)F(\alpha+1,\beta,\gamma)-(\gamma-\beta)F(\alpha,\beta-1,\gamma)
\end{array}
\]
\[
\begin{array}{rcl}
{[4]}&0=&\gamma(\alpha-(\gamma-\beta)x)F-\alpha\gamma(1-x)F(\alpha+1,\beta,\gamma)\\
&&{}+(\gamma-\alpha)(\gamma-\beta)xF(\alpha,\beta,\gamma+1)
\end{array}
\]
\[
\begin{array}{rcl}
{[5]}&0=&(\gamma-\alpha-1)F+\alpha F(\alpha+1,\beta,\gamma)-(\gamma-1)F(\alpha,\beta,\gamma-1)
\end{array}
\]
\[
\begin{array}{rcl}
{[6]}&0=&(\gamma-\alpha-\beta)F-(\gamma-\alpha)F(\alpha-1,\beta,\gamma)+\beta(1-x)F(\alpha,\beta+1,\gamma)
\end{array}
\]
\[
\begin{array}{rcl}
{[7]}&0=&(\beta-\alpha)(1-x)F-(\gamma-\alpha)F(\alpha-1,\beta,\gamma)+(\gamma-\beta)F(\alpha,\beta-1,\gamma)
\end{array}
\]
\[
\begin{array}{rcl}
{[8]}&0=&\gamma(1-x)F-\gamma F(\alpha-1,\beta,\gamma)+(\gamma-\beta)xF(\alpha,\beta,\gamma+1)
\end{array}
\]
\[
\begin{array}{rcl}
{[9]}&0=&(\alpha-1-(\gamma-\beta-1)x)F+(\gamma-\alpha)F(\alpha-1,\beta,\gamma)\\
&&{}-(\gamma-1)(1-x)F(\alpha,\beta,\gamma-1)
\end{array}
\]
\[
\begin{array}{rcl}
{[10]}&0=&(\gamma-2\beta+(\beta-\alpha)x)F+\beta(1-x)F(\alpha,\beta+1,\gamma)\\
&&{}-(\gamma-\beta)F(\alpha,\beta-1,\gamma)
\end{array}
\]
\[
\begin{array}{rcl}
{[11]}&0=&\gamma(\beta-(\gamma-\alpha)x)F-\beta\gamma(1-x)F(\alpha,\beta+1,\gamma)\\
&&{}-(\gamma-\alpha)(\gamma-\beta)F(\alpha,\beta,\gamma+1)
\end{array}
\]
\[
\begin{array}{rcl}
{[12]}&0=&(\gamma-\beta-1)F+\beta F(\alpha,\beta+1,\gamma)-(\gamma-1)F(\alpha,\beta,\gamma-1)
\end{array}
\]
\[
\begin{array}{rcl}
{[13]}&0=&\gamma(1-x)F-\gamma F(\alpha,\beta-1,\gamma)+(\gamma-\alpha)xF(\alpha,\beta,\gamma+1)
\end{array}
\]
\[
\begin{array}{rcl}
{[14]}&0=&(\beta-1-(\gamma-\alpha-1)x)F+(\gamma-\beta)F(\alpha,\beta-1,\gamma)\\
&&{}-(\gamma-1)(1-x)F(\alpha,\beta,\gamma-1)
\end{array}
\]
\[
\begin{array}{rcl}
{[15]}&0=&\gamma(\gamma-1-(2\gamma-\alpha-\beta-1)x)F\\
&&{}+(\gamma-\alpha)(\gamma-\beta)xF(\alpha,\beta,\gamma+1)\\
&&{}-\gamma(\gamma-1)(1-x)F(\alpha,\beta,\gamma-1)
\end{array}
\]
\SixthArticle{8}
Now for the proof of these formulas. Let
\[
\frac{(\alpha+1)(\alpha+2)\mathinner{.\,.\,.}(\alpha+m-1)\beta(\beta+1)\mathinner{.\,.\,.\,.\,.}(\beta+m-2)}{1\mathbin{.}2\mathbin{.}3\mathinner{.\,.\,.\,.\,.}m\mathbin{.}\gamma(\gamma+1)\mathinner{.\,.\,.\,.\,.}(\gamma+m-1)}=M
\].
\SourceLeaf{140}
Then the coefficient of the power $x^m$ is
\[
\begin{array}{lll}
\mbox{in }F&.\quad.\quad.\quad.\quad.\quad.&\alpha(\beta+m-1)M\\
\mbox{in }F(\alpha,\beta-1,\gamma)&.\quad.&\alpha(\beta-1)M\\
\mbox{in }F(\alpha+1,\beta,\gamma)&.\quad.&(\alpha+m)(\beta+m-1)M\\
\mbox{in }F(\alpha,\beta,\gamma-1)&.\quad.&\frac{\alpha(\beta+m-1)(\gamma+m-1)M}{\gamma-1}
\end{array}
\]
while the coefficient of the power $x^{m-1}$ in $F(\alpha+1,\beta,\gamma)$, or the coefficient of the power $x^m$ in $xF(\alpha+1,\beta,\gamma)$,
\[
=m(\gamma+m-1)M
\]
The truth of formulas 5 and 3 follows at once. Interchanging $\alpha$ and $\beta$ gives formula 12 from 5, and elimination between these two gives 2. In the same way, the same interchange gives 6 from 3; combining 6 and 12 gives 9, and interchange then gives 14; combining these gives 7. Finally, 2 and 6 give 1, and interchange then gives 10. Formula 8 can be derived by considering the coefficients, as formulas 5 and 3 were above (and, if desired, \textit{all} 15 formulas could be obtained in the same way), or more elegantly from those already known, as follows. Replacing the element $\alpha$ by $\alpha-1$ and $\gamma$ by $\gamma+1$ in formula 5 gives
\[
0=(\gamma-\alpha+1)F(\alpha-1,\beta,\gamma+1)+(\alpha-1)F(\alpha,\beta,\gamma+1)-\gamma F(\alpha-1,\beta,\gamma)
\]
Replacing only $\gamma$ by $\gamma+1$ in formula 9, however, gives
\[
\begin{aligned}0={}&(\alpha-1-(\gamma-\beta)x)F(\alpha,\beta,\gamma+1)\\
&{}+(\gamma-\alpha+1)F(\alpha-1,\beta,\gamma+1)-\gamma(1-x)F(\alpha,\beta,\gamma)\end{aligned}
\]
Subtracting these formulas immediately gives 8, and interchange then gives 13. From 1 and 8 comes 4, and interchange then gives 11. Finally, 8 and 9 give 15.
\SixthArticle{9}
If $\alpha'-\alpha$, $\beta'-\beta$, $\gamma'-\gamma$, as well as $\alpha''-\alpha$, $\beta''-\beta$, $\gamma''-\gamma$, are integers (positive or negative), it is possible to pass from the function $F(\alpha,\beta,\gamma)$ to the function $F(\alpha',\beta',\gamma')$, and likewise from this to the function $F(\alpha'',\beta'',\gamma'')$, through a sequence of similar functions, each contiguous to the preceding and following one. Namely, first one element, for example $\alpha$, is repeatedly changed by unity until $F(\alpha,\beta,\gamma)$ becomes $F(\alpha',\beta,\gamma)$; then the second element is changed until\SourceLeaf{141}
$F(\alpha',\beta',\gamma)$ is reached, and finally the third element is changed until $F(\alpha',\beta',\gamma')$ is reached; and likewise from this to $F(\alpha'',\beta'',\gamma'')$. Since Article 7 gives linear equations between the first, second and third functions, and in general between any three successive functions of this sequence, it is easy to see that elimination can yield a linear equation between the functions $F(\alpha,\beta,\gamma)$, $F(\alpha',\beta',\gamma')$, $F(\alpha'',\beta'',\gamma'')$. Thus, generally speaking, from two functions whose first three elements differ by integers, any other function with the same property can be derived, provided the fourth element remains the same. Here, however, it is enough for us to have established this remarkable truth in general; we do not dwell here on shortcuts by which the operations needed for this purpose can be made as brief as possible.
\SixthArticle{10}
Suppose, for example, that the functions
\[
F(\alpha,\beta,\gamma),\ F(\alpha+1,\beta+1,\gamma+1),\ F(\alpha+2,\beta+2,\gamma+2)
\]
are given, and a linear equation between them is required. Let us connect them by contiguous functions as follows:
\[
\begin{array}{rl}
F(\alpha,\beta,\gamma)&=F
\end{array}
\]
\[
\begin{array}{rl}
F(\alpha+1,\beta,\gamma)&=F'
\end{array}
\]
\[
\begin{array}{rl}
F(\alpha+1,\beta+1,\gamma)&=F''
\end{array}
\]
\[
\begin{array}{rl}
F(\alpha+1,\beta+1,\gamma+1)&=F'''
\end{array}
\]
\[
\begin{array}{rl}
F(\alpha+2,\beta+1,\gamma+1)&=F''''
\end{array}
\]
\[
\begin{array}{rl}
F(\alpha+2,\beta+2,\gamma+1)&=F'''''
\end{array}
\]
\[
\begin{array}{rl}
F(\alpha+2,\beta+2,\gamma+2)&=F''''''
\end{array}
\]
We thus have five linear equations (from formulas 6, 13, 5 of Article 7):
\[
\begin{array}{rcl}
\mathrm{I.}&0=&(\gamma-\alpha-1)F-(\gamma-\alpha-1-\beta)F'-\beta(1-x)F''
\end{array}
\]
\[
\begin{array}{rcl}
\mathrm{II.}&0=&\gamma F'-\gamma(1-x)F''-(\gamma-\alpha-1)xF'''
\end{array}
\]
\[
\begin{array}{rcl}
\mathrm{III.}&0=&\gamma F''-(\gamma-\alpha-1)F'''-(\alpha+1)F''''
\end{array}
\]
\[
\begin{array}{rcl}
\mathrm{IV.}&0=&(\gamma-\alpha-1)F'''-(\gamma-\alpha-2-\beta)F''''-(\beta+1)(1-x)F'''''
\end{array}
\]
\[
\begin{array}{rcl}
\mathrm{V.}&0=&(\gamma+1)F''''-(\gamma+1)(1-x)F'''''-(\gamma-\alpha-1)xF''''''
\end{array}
\]
Eliminating $F'$ from I and II gives

\SourceLeaf{142}
\[
\mathrm{VI.}\quad 0=\gamma F-\gamma(1-x)F''-(\gamma-\alpha-\beta-1)xF'''
\]
Eliminating $F''$ from this and III gives
\[
\mathrm{VII.}\quad 0=\gamma F-(\gamma-\alpha-1-\beta x)F'''-(\alpha+1)(1-x)F''''
\]
Further, eliminating $F'''''$ from IV and V gives
\[
\mathrm{VIII.}\quad 0=(\gamma+1)F'''-(\gamma+1)F''''+(\beta+1)xF''''''
\]
Eliminating $F''''$ from this and VII gives
\[
\begin{aligned}&\mathrm{IX.}\quad 0=\gamma(\gamma+1)F\\&{}-(\gamma+1)(\gamma-(\alpha+\beta+1)x)F'''\\&{}-(\alpha+1)(\beta+1)x(1-x)F''''''\end{aligned}
\]
